% TODO make hw stylesheet
\documentclass[10pt]{scrartcl}
\usepackage[a4paper, total={6in, 8in}]{geometry}
\usepackage[usenames,dvipsnames,svgnames]{xcolor}
\usepackage[shortlabels]{enumitem}
\usepackage[framemethod=TikZ]{mdframed}
\usepackage{amsmath,amssymb,amsthm}
\usepackage[colorlinks]{hyperref}
\usepackage{multicol}
\usepackage{tabularx}

\hypersetup{
	colorlinks=true,
	linkcolor=blue,
	filecolor=magenta,      
	urlcolor=magenta,
	pdftitle={MAT 322 HW}
}

\usepackage{microtype}
\usepackage{mathtools}
\usepackage[headsepline]{scrlayer-scrpage}
\usepackage{thmtools}
\usepackage[textsize=footnotesize]{todonotes}

\mdfdefinestyle{mdbluebox}{roundcorner=10pt,innerbottommargin=9pt,
	linecolor=blue,backgroundcolor=TealBlue!5,}
\declaretheoremstyle[headfont=\sffamily\bfseries\color{MidnightBlue},
mdframed={style=mdbluebox},]{thmbluebox}

\mdfdefinestyle{mdbloobox}{frametitlefont=\bfseries,innerbottommargin=8pt,
	nobreak=true,backgroundcolor=TealBlue!5,linecolor=blue,}
\declaretheoremstyle[headfont=\bfseries\color{MidnightBlue},
mdframed={style=mdbloobox},headpunct={\\[3pt]},postheadspace=0pt,]{thmbloobox}

\mdfdefinestyle{mdredbox}{frametitlefont=\bfseries,innerbottommargin=8pt,
	nobreak=true,backgroundcolor=Salmon!5,linecolor=RawSienna,}
\declaretheoremstyle[headfont=\bfseries\color{RawSienna},
mdframed={style=mdredbox},headpunct={\\[3pt]},postheadspace=0pt,]{thmredbox}

\mdfdefinestyle{mdgreenbox}{linecolor=ForestGreen,backgroundcolor=ForestGreen!5,
	linewidth=2pt,rightline=false,leftline=true,topline=false,bottomline=false,}
\declaretheoremstyle[headfont=\bfseries\sffamily\color{ForestGreen!70!black},
mdframed={style=mdgreenbox},headpunct={ --- },]{thmgreenbox}

\mdfdefinestyle{mdorangebox}{linecolor=RedOrange,backgroundcolor=RedOrange!5,
	linewidth=2pt,rightline=false,leftline=true,topline=false,bottomline=false,}
\declaretheoremstyle[headfont=\bfseries\sffamily\color{RedOrange!70!black},
mdframed={style=mdorangebox},headpunct={ --- },]{thmorangebox}

\mdfdefinestyle{mdblackbox}{linecolor=black,backgroundcolor=RedViolet!5!gray!5,
	linewidth=3pt,rightline=false,leftline=true,topline=false,bottomline=false,}
\declaretheoremstyle[mdframed={style=mdblackbox}]{thmblackbox}

\declaretheoremstyle[spaceabove=6pt,spacebelow=6pt]{basehead}

\declaretheorem[style=basehead,name=Problem]{problem}
\declaretheorem[style=thmredbox,name=Problem,numbered=no]{problem*}
\declaretheorem[style=thmbloobox,name=Setup]{setup} % for config geo
\declaretheorem[style=thmredbox,name=Example,sibling=problem]{example}
\declaretheorem[style=thmredbox,name=$\bigstar$ Problem,sibling=problem]{niceprob}
\declaretheorem[style=thmbluebox,name=Theorem,sibling=problem]{theorem}
\declaretheorem[style=thmbluebox,name=Corollary,sibling=theorem]{corollary}
\declaretheorem[style=thmbluebox,name=Proposition,sibling=theorem]{prop}
\declaretheorem[style=thmbluebox,name=Lemma,numberwithin=problem]{lemma}
\declaretheorem[style=thmbluebox,name=Theorem,numbered=no]{theorem*}
\declaretheorem[style=thmbluebox,name=Lemma,numbered=no]{lemma*}
\declaretheorem[style=thmgreenbox,name=Claim,numberwithin=problem]{claim}
\declaretheorem[style=thmgreenbox,name=Claim,numbered=no]{claim*}
\declaretheorem[style=thmblackbox,name=Remark,numberwithin=problem]{remark}
\declaretheorem[style=thmblackbox,name=Remark,numbered=no]{remark*}
\declaretheorem[style=thmgreenbox,name=Definition,sibling=theorem]{definition}
\declaretheorem[style=thmgreenbox,name=Definition,numbered=no]{definition*}
\declaretheorem[style=thmorangebox,name=Warning,sibling=theorem]{warning}
\declaretheorem[style=thmorangebox,name=Warning,numbered=no]{warning*}
\declaretheorem[style=thmorangebox,name=Note,numbered=no]{note}
\declaretheorem[style=thmblackbox,name=Exercise,sibling=problem]{exercise}
\declaretheorem[style=thmblackbox,name=Exercise,numbered=no]{exercise*}
\declaretheorem[style=thmblackbox,name=Question,sibling=problem]{question}
\declaretheorem[style=thmblackbox,name=Question,numbered=no]{question*}

\addtolength{\textheight}{3.14cm}
\providecommand{\ol}{\overline}
\providecommand{\ul}{\underline}
\providecommand{\eps}{\varepsilon}
\providecommand{\half}{\frac{1}{2}}
\providecommand{\dang}{\measuredangle} %% Directed angle
\providecommand{\CC}{\mathbb C}
\providecommand{\FF}{\mathbb F}
\providecommand{\NN}{\mathbb N}
\providecommand{\QQ}{\mathbb Q}
\providecommand{\RR}{\mathbb R}
\providecommand{\ZZ}{\mathbb Z}
\providecommand{\dg}{^\circ}
\providecommand{\ii}{\item}
\providecommand{\setdiff}{\smallsetminus}
\providecommand{\alert}[1]{{\sffamily\textbf{\textcolor{blue}{#1}}}}
\providecommand{\inv}{^{-1}}
\providecommand{\mb}[1]{\mathbf{#1}}
\providecommand{\dif}{\;d}

\reversemarginpar

\newenvironment{soln}{\begin{proof}[Solution]}{\end{proof}}
\newenvironment{walkthrough}{\noindent \textbf{Walkthrough.}}{\medskip}
\newlist{walk}{enumerate}{3}
\setlist[walk]{label=\bfseries (\alph*)}

\providecommand{\pow}{\operatorname{Pow}}
\providecommand{\vocab}[1]{{\textbf{\textcolor{ForestGreen}{#1}}}}
\providecommand{\lcm}{\operatorname{lcm}}
\providecommand{\abs}[1]{\left|#1\right|}

\usepackage{asymptote}
\begin{asydef}
	defaultpen(fontsize(10pt));
	unitsize(1cm);
	size(8cm); // set a reasonable default
	usepackage("amsmath");
	usepackage("amssymb");
	settings.tex="pdflatex";
	settings.outformat="pdf";
	// Replacement for olympiad+cse5 which is not standard
	import geometry;
	// recalibrate fill and filldraw for conics
	void filldraw(picture pic = currentpicture, conic g, pen fillpen=defaultpen, pen drawpen=defaultpen)
	{ filldraw(pic, (path) g, fillpen, drawpen); }
	void fill(picture pic = currentpicture, conic g, pen p=defaultpen)
	{ filldraw(pic, (path) g, p); }
	// some geometry
	pair foot(pair P, pair A, pair B) { return foot(triangle(A,B,P).VC); }
	pair orthocenter(pair A, pair B, pair C) { return orthocentercenter(A,B,C); }
	pair centroid(pair A, pair B, pair C) { return (A+B+C)/3; }
	// cse5 abbreviations
	path CP(pair P, pair A) { return circle(P, abs(A-P)); }
	path CR(pair P, real r) { return circle(P, r); }
	pair IP(path p, path q) { return intersectionpoints(p,q)[0]; }
	pair OP(path p, path q) { return intersectionpoints(p,q)[1]; }
	path Line(pair A, pair B, real a=0.6, real b=a) { return (a*(A-B)+A)--(b*(B-A)+B); }
	// cse5 more useful functions
	picture CC() {
		picture p=rotate(0)*currentpicture;
		currentpicture.erase();
		return p;
	}
	pair MP(Label s, pair A, pair B = plain.S, pen p = defaultpen) {
		Label L = s;
		L.s = "$"+s.s+"$";
		label(L, A, B, p);
		return A;
	}
	pair Drawing(Label s = "", pair A, pair B = plain.S, pen p = defaultpen) {
		dot(MP(s, A, B, p), p);
		return A;
	}
	path Drawing(path g, pen p = defaultpen, arrowbar ar = None) {
		draw(g, p, ar);
		return g;
	}
\end{asydef}

\usepackage{mathrsfs}

\ihead{\footnotesize MAT 322 HW03}
\ohead{\footnotesize Last compiled \today}

\title{\vspace{-2em}MAT 322 HW03}
\author{\large Aditya Pahuja}
\date{\large Due 02/18}
\begin{document}
\maketitle

\begin{problem}
	Let $f\colon\RR^2\to\RR^2$ be defined by the equation
	\[f(x,y) = (x^2 - y^2, 2xy).\]
	\begin{enumerate}[(a)]
		\ii Show that $f$ is injective on the set $U$ consisting of
		all $(x,y)$ with $x>0$.
		\ii What is the set $V = f(U)$?
		\ii If $g$ is the inverse function, find $(\mathrm dg)_{(0,1)}$.
	\end{enumerate}
\end{problem}

\begin{soln}
	(a)	Observe that we can write $f$ as the sequence of maps
	\[(x, y)\mapsto x + yi\mapsto (x + yi)^2 = (x^2 - y^2) + i(2xy)
	\mapsto (x^2 - y^2, 2xy).\]
	The first and last two maps, which swap between representing a point
	as a pair of coordinates and as a complex number,
	are bijective maps $U\to U$: we can clearly see that $(u, v)\mapsto u+vi$
	is invertible with $u+vi\mapsto(u,v)$ as its inverse.
	The map in the middle, $z\mapsto z^2$, is injective on $U$
	because for any $a\in\CC$, $z^2 = a$
	has exactly two solutions $z_1,z_2\in\CC$
	that are negations of one another,
	so at most one of them has a positive real part.
	Compositions of injective maps are injective,
	so $f$ is injective on $U$.
	
	(b) $V$ is the set
	\[\RR^2\setdiff\{(r,0)\in\RR^2 : r \le 0\}\]
	Indeed, for $a\in\CC$, the only time that $z^2 = a$
	does not have a solution in $\CC$ with positive real part
	is when both solutions have a real part of zero,
	which happens if and only if $a$ is a nonpositive real number.
	
	(c) Letting $I$ be the $2\times 2$ identity matrix, we have
	\[I = (\mathrm d(f\circ g))_{(0,1)} =
	(\mathrm df)_{(g(0,1))}\circ(\mathrm dg)_{(0,1)}
	= (\mathrm df)_{(1,1)}\circ(\mathrm dg)_{(0,1)},\]
	so $(\mathrm dg)_{(0,1)} = [(\mathrm df)_{(1,1)}]\inv$.
	The derivative of $f$ in general is the matrix
	\[\begin{pmatrix}
		2x&2y\\
		-2y&2x
	\end{pmatrix},\]
	so the answer is
	\[(\mathrm dg)_{(0,1)} = [(\mathrm df)_{(1,1)}]\inv = \begin{pmatrix}
		2&2\\-2&2
	\end{pmatrix}\inv = \boxed{\begin{pmatrix}
		1/4&1/4\\-1/4&1/4
	\end{pmatrix}}.\]
\end{soln}

\begin{problem}
	Let $f\colon\RR^2\to\RR^2$ be defined by the equation
	\[f(x,y) = (e^x\cos y, e^x\sin y).\]
	\begin{enumerate}[(a)]
		\ii Show that $f$ is injective on the set $U$ consisting of
		all $(x,y)$ with $0<y<2\pi$.
		\ii What is the set $V = f(U)$?
		\ii If $g$ is the inverse function, find $(\mathrm dg)_{(0,1)}$.
	\end{enumerate}
\end{problem}

\begin{soln}
	(a) Suppose that $f(x,y) = f(x',y')$, so that
	\[(e^x\cos y, e^x\sin y) = (e^{x'}\cos y', e^{x'}\sin y').\]
	Then the magnitudes of these vectors are the same, so
	$e^x = e^{x'}$, which implies $x=x'$ because exponential functions
	are injective on $\RR$. This in turn implies that
	$(\cos y,\sin y) = (\cos y', \sin y')$, which tells us that
	$y = y'$ in light of the restriction on the size of $y$.
	
	(b) $V$ is the set $\RR^2\setdiff\{(t,0):t\ge0\}$. Indeed, for any
	point $(r\cos\theta, r\sin\theta)$ with $r > 0$ and $\theta\ne 0$
	(hence in $V$), we can set $x = \ln r$, $y = \theta$ to find a point
	whose image under $f$ is in $V$. Conversely, if $(a, b)$ is not in $V$,
	then we either get $(a, b) = (0,0)$, so $e^x = 0$,
	which has no real solutions, or $(a, b) = (t, 0)$ for some $t > 0$, giving
	\[\cos y = 1,\quad \sin y = 0\]
	so $y = 2\pi n$ for some $n\in\ZZ$, which is not in the interval $(0,2\pi)$.
	
	(c) Noting that $g(0,1) = (0, \frac\pi2)$, we get
	\[(\mathrm dg)_{(0,1)} = [(\mathrm df)_{(0,\frac\pi2)}]\inv =
	\begin{pmatrix}
		e^0\cos\pi/2&e^0\sin\pi/2\\
		-e^0\sin\pi/2&e^0\cos\pi/2
	\end{pmatrix}\inv =\begin{pmatrix}
		0&1\\-1&0
	\end{pmatrix}\inv = \boxed{\begin{pmatrix}
		0&-1\\1&0
	\end{pmatrix}}.\qedhere\]
\end{soln}

\begin{problem}
	Let $f\colon\RR^n\to\RR^n$ be given by the equation $f(x) = \|x\|^2\cdot x$.
	Show that $f$ is continuously differentiable and that
	$f$ carries the open ball $B_1(0)$ onto itself in an injective fashion.
	Show, however, that the inverse of $f$ is not differentiable at $0$.
\end{problem}

\begin{soln}
	Letting $x = (x_1, \dots, x_n)$, we can write down
	\[f(x) = \begin{pmatrix}
		x_1(x_1^2 + x_2^2 + \cdots + x_n^2)\\
		x_2(x_1^2 + x_2^2 + \cdots + x_n^2)\\
		\vdots\\
		x_n(x_1^2 + x_2^2 + \cdots + x_n^2)\\
	\end{pmatrix}.\]
	Then we can compute the Jacobian matrix
	\[(\mathrm df)_x = \begin{pmatrix}
		3x_1^2 & 2x_1x_2 & \cdots & 2x_1x_n\\
		2x_2x_1 & 3x_2^2 & \cdots & 2x_2x_n\\
		\vdots & \vdots & \ddots & \vdots\\
		2x_nx_1 & 2x_nx_2 & \cdots & 3x_n^2\\
	\end{pmatrix}\]
	whose entries are clearly continuous in $x_1$, \dots, $x_n$,
	so $f$ is continuously differentiable.
	
	To show that $f$ is a bijection on $B_1(0)$, I claim that
	the function $g(x) = \frac{x}{\|x\|^{2/3}}$ is the inverse of $f$
	on this ball. Indeed,
	\begin{align*}
		g(f(x)) &= g(\|x\|^2\cdot x) =
		\frac{\|x\|^2\cdot x}{\|\|x\|^2\cdot x\|^{2/3}} = x\\
		f(g(x)) &= f\left(\frac{x}{\|x\|^{2/3}}\right) =
		\left\|\frac{x}{\|x\|^{2/3}}\right\|^2\cdot \frac{x}{\|x\|^{2/3}} = x.
	\end{align*}
	
	This inverse is not differentiable at $0$ because its partials
	are not defined at zero. For example, the limit
	\[\lim_{h\to 0}\frac{f(h, 0, \dots, 0) - f(0, 0, \dots, 0)}{h}
	= \lim_{h\to 0} (1/h^{2/3}, 0, \dots, 0)\]
	does not exist.
\end{soln}

\begin{problem}
	Let $g\colon\RR^2\to\RR^2$ be given by the equation
	\[g(x,y) = (2ye^{2x}, xe^y)\]
	and let $f\colon\RR^2\to\RR^3$ be given by the equation
	\[f(x,y) = (3x-y^2, 2x+y, xy+y^3).\]
	\begin{enumerate}[(a)]
		\ii Show that there is a neighborhood of $(0,1)$ that $g$ carries
		injectively onto a neighborhood of $(2,0)$.
		\ii Find $(\mathrm d(f\circ g\inv))_{(2,0)}$.
	\end{enumerate}
\end{problem}

\begin{soln}
	Observe that $g$ is a class $C^1$ function and that
	\[(\mathrm dg)_{(0,1)} = \begin{pmatrix}
		4&2\\e&0
	\end{pmatrix}\]
	is invertible, so the inverse function theorem says that
	there is a neighborhood $U$ of $(0,1)$ that $g$ carries bijectively
	to a neighborhood of $g(0,1) = (2,0)$ (which is part (a))
	and that $g$ has a continuously differentiable inverse on the set $g(U)$.
	
	For part (b) we can compute
	\begin{align*}
		(\mathrm d(f\circ g\inv))_{(2,0)}
		&= (\mathrm df)_{g\inv(2,0)}\circ (\mathrm dg\inv)_{(2,0)}\\
		&= (\mathrm df)_{g\inv(2,0)}\circ [(\mathrm dg)_{g\inv(2,0)}]\inv\\
		&= (\mathrm df)_{(0,1)}\circ [(\mathrm dg)_{(0,1)}]\inv\\
		&= \begin{pmatrix}
			3&-2\\2&1\\1&2
		\end{pmatrix}\begin{pmatrix}
			0&1/e\\1/2&-2/e
		\end{pmatrix}\\
		&= \begin{pmatrix}
			-1&7/e\\1/2&0\\1&-3/e
		\end{pmatrix}.\qedhere
	\end{align*}
\end{soln}

\begin{problem}
	Let $U\subseteq\RR^n$ be open and let $f\colon U\to\RR^n$
	be continuously differentiable. If $(\mathrm df)_x$ is invertible
	for all $x\in U$, prove that $f(U)$ is an open set.
\end{problem}

\begin{soln}
	Let $x$ be a point in $U$.
	Since $f$ is a $C^1$ function with invertible derivative at $x$,
	we can see from the proof of the inverse function theorem that
	there is an open ball $B_x\subseteq U$ such that
	$f(B_x)\subseteq f(U)$ is open. This means that we can write
	\[f(U) = \bigcup_{x\in U} f(B_x)\]
	as a union of open sets, hence $f(U)$ is also open.
\end{soln}


























\end{document}